\documentclass[a4paper,11pt]{article}

\usepackage[T1]{fontenc}
\usepackage[utf8]{inputenc}
\usepackage{lmodern}
\usepackage{microtype}
\usepackage{geometry}
% \geometry{margin=2.5cm}
\usepackage{hyperref}
\usepackage{enumitem}
\setlist{nosep}

\title{IncludeOS Kernel Startup Sequence\\(x86\_64, Multiboot, up to APIC Timer Calibration)}
\author{Mazunki Hoksaas}
\date{December 4, 2025}

\begin{document}
\maketitle

\section{Scope and Intent}

This document reconstructs the IncludeOS kernel startup sequence on
the 64-bit x86 architecture (x86\_64) for a Multiboot boot scenario,
based on:

\begin{itemize}
  \item a GDB (GNU Debugger) trace with breakpoints across early kernel
        code, and
  \item the serial boot log printed by the kernel during boot.
\end{itemize}

IncludeOS also has early support and stub code for other architectures
(e.g.\ AArch64), but the trace and log I am working from exercise only
the \texttt{x86\_pc} platform with the x86\_64 architecture. Everything
here should therefore be read as describing \emph{the x86\_64 Multiboot
configuration specifically}. Where relevant, I point out that other
architectures have different entry stubs and platform glue.

The goal is not just to list which functions run in which order, but to
explain \emph{why} the IncludeOS kernel is structured this way, and what
each phase guarantees for later layers such as the C standard library
(\texttt{libc}) and the C++ standard library (\texttt{libc++}). This is
intended both as IncludeOS documentation and as material for my master
thesis, so I explicitly discuss design choices and ordering, not only
mechanics.

Where the exact behaviour is unclear from logs and a brief reading of
the code at commit
\texttt{62337fceb8f2aa72acd90de07ef9a3f69ec98197}, I leave small
\texttt{TODO} notes as anchors for further work.

The current logs end while the kernel is still calibrating the APIC
(Advanced Programmable Interrupt Controller) timer. The system is
clearly alive and servicing xAPIC (x86 Advanced Programmable Interrupt
Controller) interrupts, but has not yet reached a fully booted state (no
service entry or main event loop is visible in this trace). Later phases
are therefore intentionally left incomplete until I have extended the
trace or annotated the corresponding source in more detail.

\section{From Firmware to \texorpdfstring{\texttt{kernel\_start}}{kernel\_start}}

\subsection{Bootloader and Multiboot handoff}

On the \texttt{x86\_pc} platform, the IncludeOS kernel image is loaded by
a Multiboot-capable loader (either GRUB or the IncludeOS chainloader).
The loader:

\begin{itemize}
  \item loads the IncludeOS ELF image into memory,
  \item optionally loads additional modules (e.g.\ the service binary),
  \item constructs the Multiboot information structure, and
  \item jumps into the entry point declared in the Multiboot header.
\end{itemize}

The Multiboot header for this platform is defined in
\texttt{src/platform/x86\_pc/start.asm} in a dedicated
\texttt{.multiboot} section. It contains:

\begin{itemize}
  \item the Multiboot magic (\texttt{0x1BADB002}),
  \item flags (requesting memory information and alignment),
  \item pointers to the ELF load range (\texttt{\_LOAD\_START\_},
        \texttt{\_LOAD\_END\_}) and the end-of-image symbol
        (\texttt{\_end}),
  \item the entry function \texttt{\_start},
  \item and a secondary entry \texttt{fast\_kernel\_start} used for live
        updates.
\end{itemize}

When the loader jumps into the image, the CPU is in 32-bit protected
mode and the Multiboot convention guarantees:

\begin{itemize}
  \item \texttt{EAX} contains the Multiboot magic (\texttt{0x2BADB002}),
  \item \texttt{EBX} contains a pointer to the Multiboot information
        structure.
\end{itemize}

My serial log reflects this initial handover, for example:

\begin{verbatim}
* Multiboot begin: 0x9500
* Multiboot cmdline @ 0x003c6034: /nix/store/.../chainloader ""
* Module @ 0x3c7000
  * Args: ./result/hello_includeos.elf.bin ""
  * End: 0x17997d0
* ELF syms header CRC OK
< multiboot >Booted with multiboot
\end{verbatim}

At this point, none of the 64-bit kernel code has run yet; execution is
still in the 32-bit platform-specific entry stub in
\texttt{start.asm}.

\subsection{32-bit entry stub on \texorpdfstring{\texttt{x86\_pc}}{x86\_pc}}

The 32-bit boot stub for x86 is implemented in
\texttt{src/platform/x86\_pc/start.asm}. Its responsibilities are:

\begin{enumerate}
  \item Install a minimal 32-bit Global Descriptor Table (GDT).
  \item Set up a 32-bit stack in a safe low-memory region.
  \item Enable floating-point and SIMD support (FPU, SSE, XSAVE, AVX).
  \item Preserve the Multiboot parameters (\texttt{EAX}, \texttt{EBX}).
  \item Transfer control to the architecture-specific long-mode setup
        routine \texttt{\_\_arch\_start}.
\end{enumerate}

The stub builds a small GDT consisting of:

\begin{itemize}
  \item a null descriptor,
  \item a flat 32-bit code segment (selector \texttt{0x08}),
  \item a flat 32-bit data segment (selector \texttt{0x10}),
  \item a data segment (selector \texttt{0x18}) later repurposed as
        \texttt{GS} for local CPU data.
\end{itemize}

It loads this GDT via \texttt{lgdt [gdtr]} and then reloads all segment
registers. The stack pointer is set to \texttt{0x9D3F0}, just below the
EBDA (Extended BIOS Data Area). The comments in \texttt{start.asm}
explicitly note that Multiboot may use \texttt{0x9D400--0x9FFFF}, so the
stack is placed directly below that to avoid corrupting Multiboot data.

Before calling into any C or C++ code, the stub enables floating-point
and SIMD support through CR0/CR4 and checks CPUID flags to optionally
enable XSAVE and AVX. It also sets global flags like \texttt{\_\_xsave\_enabled}
and \texttt{\_\_avx\_enabled} based on what the CPU reports.

Finally, it stores the Multiboot parameters into globals:

\begin{verbatim}
__multiboot_magic = EAX
__multiboot_addr  = EBX
\end{verbatim}

and calls the architecture-specific entry \texttt{\_\_arch\_start}. If
\texttt{\_\_arch\_start} ever returns, the stub prints a panic message
over the serial port and halts.

\subsection{Transition to 64-bit long mode on x86\_64}

The function \texttt{\_\_arch\_start} is defined in
\texttt{src/arch/x86\_64/arch\_start.asm}. Its job is to transition from
32-bit protected mode into 64-bit long mode with a controlled paging and
segmentation setup, and then call \texttt{kernel\_start} in C++.

The sequence in \texttt{arch\_start.asm} is:

\begin{enumerate}
  \item Disable any existing paging by clearing the PG bit in CR0.
  \item Construct a minimal long-mode page-table hierarchy:
    \begin{itemize}
      \item a PML4 (Page Map Level 4) at physical address
            \texttt{0x1000},
      \item a PDPT (Page Directory Pointer Table) at \texttt{0x2000},
      \item a range of Page Directories starting at \texttt{0x100000},
            each filled with 2~MiB ``huge'' pages.
    \end{itemize}
    This creates a large identity-mapped region so that virtual and
    physical addresses match for the early kernel.
  \item Enable Physical Address Extension (PAE) in CR4 and program the
        IA32\_EFER (Extended Feature Enable Register) MSR to turn on
        long mode and no-execute (NX).
  \item Re-enable paging and supervisor write-protect via CR0.
  \item Load a 64-bit GDT with code and data descriptors and perform a
        far jump to the \texttt{long\_mode} label.
\end{enumerate}

In the \texttt{long\_mode} section, the code:

\begin{itemize}
  \item disables interrupts (\texttt{cli}),
  \item reloads all segment registers (\texttt{ds}, \texttt{es},
        \texttt{fs}, \texttt{gs}, \texttt{ss}) with the 64-bit data
        segment,
  \item installs a new 64-bit stack by setting \texttt{RSP} to the
        linker-provided symbol \texttt{\_ELF\_START\_} (top of the loaded
        ELF image) and \texttt{RBP} to match,
  \item programs IA32\_STAR to configure the SYSCALL/SYSRET segment
        selectors,
  \item sets IA32\_FS\_BASE to a small \texttt{tls\_table} and
        IA32\_GS\_BASE to a \texttt{smp\_table} to provide minimal
        thread-local and per-CPU storage for the kernel itself.
\end{itemize}

Finally, it loads the saved Multiboot values and calls:

\begin{verbatim}
kernel_start(__multiboot_magic, __multiboot_addr)
\end{verbatim}

using the standard System~V AMD64 calling convention (\texttt{EDI},
\texttt{ESI}).

There is also a secondary entry \texttt{fast\_kernel\_start} in the same
file, used by the live-update mechanism: it realigns the stack, forwards
two arguments (\texttt{EAX}, \texttt{EBX}) to \texttt{kernel\_start}, and
halts once \texttt{kernel\_start} returns.

\subsection{Non-x86 architectures}

IncludeOS has an \texttt{aarch64} port with its own entry stubs in:

\begin{itemize}
  \item \texttt{src/arch/aarch64/arch\_start.asm},
  \item \texttt{src/platform/aarch64\_vm/start.asm},
  \item \texttt{src/platform/aarch64\_vm/start\_aarch64.asm}.
\end{itemize}

These serve the same high-level purpose: move from the bootloader
environment into a well-defined 64-bit state and then call
\texttt{kernel\_start}. The details differ substantially (different
registers, page-table layout, exception model, etc.), and I have not yet
walked through them line by line.

\subsection{Why IncludeOS does it this way}

At this commit, x86\_64 is the primary architecture IncludeOS targets.
The kernel image itself is 64-bit C++, but the Multiboot specification
uses a 32-bit entry protocol. The architecture-specific assembly is
therefore kept as small as possible: once it has received
\texttt{magic} and \texttt{addr} from the loader and entered long mode,
it immediately hands over control to C++.

Other operating systems such as Linux or the BSD family often keep a lot
of early hardware setup in assembly: page tables, hardware probing,
stack management, etc. IncludeOS pushes almost all such logic into C++
instead. The trade-off is:

\begin{itemize}
  \item more reliance on a cooperative boot environment (virtualised
        hardware, Multiboot loader, no legacy 16-bit BIOS calls),
  \item but much better debuggability and maintainability of early boot
        logic, because it lives in regular C++.
\end{itemize}

The ARM/AArch64 path has its own assembly to match the architecture, but
is structured around the same goal: do the minimal amount of work in
assembly and defer as much as possible to C++.

\subsection{libc / \texorpdfstring{\texttt{libc++}}{libc++} assumptions at this point}

At entry to \texttt{kernel\_start}, the following are \emph{not} yet true:

\begin{itemize}
  \item there is no C heap,
  \item thread-local storage (TLS) for user-space is not configured,
  \item the final GDT (Global Descriptor Table) and IDT (Interrupt
        Descriptor Table) are not fully configured,
  \item no \texttt{libc} or \texttt{libc++} initialisation has run.
\end{itemize}

The only guarantees are:

\begin{itemize}
  \item the CPU runs in 64-bit long mode,
  \item a valid stack exists for CPU~0 (set by \texttt{arch\_start.asm}),
  \item \texttt{magic} and \texttt{addr} point to sane Multiboot data.
\end{itemize}

Everything else must be established explicitly in the subsequent phases.

\section{\texorpdfstring{\texttt{kernel\_start}}{kernel\_start}: Multiboot and Early Sanity}

\subsection{High-level flow in \texttt{kernel\_start}}

On \texttt{x86\_pc}, the C++ entry point is implemented in
\texttt{src/platform/x86\_pc/kernel\_start.cpp}:

\begin{verbatim}
extern "C"
__attribute__((no_sanitize("all")))
void kernel_start(uint32_t magic, uint32_t addr)
{
  __init_sanity_checks();
  uintptr_t free_mem_begin = reinterpret_cast<uintptr_t>(&_end);
  uintptr_t memory_end     = kernel::memory_end();

  if (magic == MULTIBOOT_BOOTLOADER_MAGIC) {
    free_mem_begin = _multiboot_free_begin(addr);
    memory_end     = _multiboot_memory_end(addr);
  } else if (kernel::is_softreset_magic(magic)) {
    memory_end = kernel::softreset_memory_end(addr);
  }

  free_mem_begin += _move_symbols(free_mem_begin);
  _init_bss();
  kernel::diag::hook<kernel::diag::post_bss>();

  size_t memsize = memory_end - free_mem_begin;
  __machine = os::Machine::create((void*)free_mem_begin, memsize);

  _init_elf_parser();
  __machine->init();
  kernel::diag::hook<kernel::diag::post_machine_init>();

  RNG::init();
  __init_crash_contexts();
  x86::idt_initialize_for_cpu(0);

  x86::init_libc(magic, addr);
}
\end{verbatim}

The rest of this section and the next few sections unpack that flow
step-by-step.

\subsection{Initial sanity and Multiboot inspection}

Inside \texttt{kernel\_start}, the kernel immediately sanity-checks the
Multiboot environment and records the available physical memory. This is
partly in \texttt{kernel\_start.cpp} and partly in
\texttt{src/kernel/multiboot.cpp}.

The sequence visible in the trace and logs is:

\begin{verbatim}
kernel_start(magic, addr)
  → __init_sanity_checks()
  → kernel::memory_end() / _end
  → _multiboot_free_begin(addr)
  → _multiboot_memory_end(addr)
  → bootinfo(addr)
  → kprintf / serial1: print Multiboot info
\end{verbatim}

The same Multiboot data is printed twice in my environment: once by the
chainloader, and once after the kernel has taken over paging and
reconstructed its own view of the memory map:

\begin{verbatim}
< multiboot >Booted with multiboot
  * Boot flags: 0x24f
  * Valid memory (%i KiB):
    0x00000000 - 0x0009fbff (639 KiB)
    0x00100000 - 0x07fdffff (129920 KiB)
  * Multiboot provided memory map  (7 entries @ 0x9000)
    0x0        - 0x9fbff    FREE
    0x9fc00    - 0x9ffff    RESERVED
    0xf0000    - 0xfffff    RESERVED
    0x100000   - 0x7fdffff  FREE
    0x7fe0000  - 0x7ffffff  RESERVED
    0xfffc0000 - 0xffffffff RESERVED
    0xfd00000000 - 0xffffffffff RESERVED
\end{verbatim}

\subsection{Why this comes before everything else}

There are two core reasons to do Multiboot parsing this early:

\begin{enumerate}
  \item \textbf{Memory safety and correctness.}  
    Every later decision about heap placement, paging, and device
    mapping depends on knowing exactly which physical ranges are usable.
    If this step is wrong or skipped, later code might:
    \begin{itemize}
      \item map over firmware memory,
      \item clobber the Multiboot info itself, or
      \item overlap modules with the heap.
    \end{itemize}

  \item \textbf{Debuggability.}  
    By printing the memory map twice (once from the chainloader and once
    from the kernel’s own context) it becomes much easier to spot
    discrepancies caused by relocation, symbol moves, or bugs in early
    memory calculations.
\end{enumerate}

In a unikernel context, there is no “initrd versus root filesystem
versus userland” to fallback on if early memory setup is wrong. The
service and the kernel share the same address space, and the heap that
\texttt{libc} uses is also the heap that the kernel uses. That pushes
correctness pressure into this phase.

\subsection{libc / \texorpdfstring{\texttt{libc++}}{libc++} implications}

\texttt{libc} and \texttt{libc++} are still completely uninitialised at this
point, but this phase defines the \emph{constraints} under which they
will later run:

\begin{itemize}
  \item the maximum amount of RAM that \texttt{brk}/\texttt{mmap} can ever
        manage,
  \item the physical regions that can be safely used for page tables and
        book-keeping,
  \item the split between “kernel reserved” and “general” memory.
\end{itemize}

From the perspective of later code, this is the moment where “the total
RAM in the VM” becomes a concrete \emph{policy} decision about how much
of that RAM may become heap.

\section{ELF Symbol Relocation and \texorpdfstring{\texttt{.bss}}{.bss} Safety}

\subsection{What happens}

Before any dynamic memory structures are built, the kernel relocates its
ELF symbols and clears the \texttt{.bss} (BSS: Block Started by Symbol)
segment.

This is handled in \texttt{src/kernel/elf.cpp} and is called from
\texttt{kernel\_start.cpp}. In the trace this appears as:

\begin{verbatim}
→ _move_symbols(free_mem_begin)
    → _get_elf_section_datasize(...)
    → _move_elf_syms_location(old, new)
    → _get_elf_section_datasize(new)   // sanity re-check

→ _init_bss()
→ kernel::diag::hook<kernel::diag::post_bss>()
\end{verbatim}

The relocation code:

\begin{itemize}
  \item inspects the ELF section headers,
  \item allocates a new region for the symbol and string tables at
        the current \texttt{free\_mem\_begin},
  \item copies the symbol and string tables to this region, and
  \item verifies that the calculated size is still consistent afterwards.
\end{itemize}

\texttt{\_init\_bss} then walks the \texttt{.bss} region and zeros it out,
bringing all implicitly-zeroed globals into a well-defined state. The
\texttt{post\_bss} diagnostic hook becomes a convenient place for sanity
checks that rely on global state but still do not require a heap.

\subsection{Why this ordering?}

There are a few subtle ordering constraints here:

\begin{enumerate}
  \item \textbf{Symbols must move before heap setup.}  
    The early memory range returned from Multiboot often includes the
    ELF image and its sections. If the heap is simply placed on “all
    free memory” without moving symbol sections, a later allocation can
    overwrite the symbol table. For a unikernel this is particularly
    painful because there is no user–kernel split where debug helpers
    might survive.

  \item \textbf{\texttt{.bss} must be zeroed before any C++ global or libc code
    runs.}  
    C and C++ treat static-storage-duration variables with no explicit
    initializer as zero-initialized. Until \texttt{\_init\_bss} runs, that
    guarantee is not actually true. Any function that reads or writes
    such globals before this is a bug.

  \item \textbf{Diagnostics after \texttt{.bss}.}  
    Running \texttt{kernel::diag::post\_bss} immediately after \texttt{.bss}
    is cleared allows debugging hooks to run in a context where:
    \begin{itemize}
      \item stacks are valid,
      \item global variables are consistent,
      \item but the heap and \texttt{libc} are still guaranteed not to be
            used.
    \end{itemize}
\end{enumerate}

\subsection{libc / \texorpdfstring{\texttt{libc++}}{libc++} requirements}

This phase establishes invariants needed by both musl and \texttt{libc++}:

\begin{itemize}
  \item \textbf{Zeroed globals.}  
    Internals such as \texttt{errno}, mutexes, and various internal
    structures in musl often live in \texttt{.bss}.

  \item \textbf{Stable symbol tables.}  
    While neither musl nor \texttt{libc++} \emph{require} symbol tables
    to function, having them preserved enables:
    \begin{itemize}
      \item meaningful crash reports and backtraces,
      \item debugging tools layered on top of IncludeOS.
    \end{itemize}
\end{itemize}

For the thesis angle: this is a concrete example of how a unikernel can
remain debuggable even though it is “just a library” linked into a
service binary. The kernel contradicts the typical “stripped release
binary” model and instead carefully preserves exactly the symbol data it
needs.

\section{Machine Construction and the Lstack Allocator}

\subsection{What happens}

With Multiboot parsed and core invariants established, the kernel now
constructs its internal representation of “the machine” and sets up an
extremely simple allocator over the physical memory region.

The relevant code lives in \texttt{src/hal/machine.hpp} (for the
interface) and the platform-specific machine implementation (not all of
which is in the bundle I used here). From the trace and logs we see
something like:

\begin{verbatim}
→ os::Machine::create(mem, size = 128839680)
    → os::Machine::Machine(this, mem, size = ...)
    → os::detail::Machine::Machine(this, raw_mem, size = ...)
        → INFO("[x86_64 PC] constructor")

→ os::Machine::Memory::Lstack()
    → util::alloc::Lstack::Lstack(mem_begin, size = 128839488)
    → donate(ptr, size)
        → bits::is_aligned(A = 64, ptr = 0x502000)
        → bits::is_aligned(A = 64, size = 128839488)
    → push(ptr, size)
\end{verbatim}

At this point, \texttt{os::Machine} has:

\begin{itemize}
  \item recorded the raw physical memory range provided by Multiboot and
        \texttt{kernel::memory\_end()},
  \item wrapped it in a small LIFO-like allocator (\texttt{Lstack}) that
        manages one or more large contiguous segments,
  \item verified that the base pointer \texttt{0x502000} and the size are
        aligned to the requested granularity.
\end{itemize}

The same base address appears later in the printed virtual memory map:

\begin{verbatim}
* Assigning heap 0x502000 -> 0x6018fff
...
0x0000502000 - 0x0006018fff,  91.090_MiB Dynamic memory
\end{verbatim}

\subsection{Why an LIFO “Lstack” allocator here?}

The choice of an Lstack allocator at this stage is deliberate:

\begin{enumerate}
  \item \textbf{Boot-time simplicity.}  
    Early boot allocators are easier to get right if they support only
    very constrained operations. The Lstack in
    \texttt{api/util/alloc\_lstack.hpp} supports a small set of
    operations (donate a big chunk of memory, hand out large chunks, and
    occasionally reclaim them) without the complexity of a general
    \texttt{malloc}.

  \item \textbf{Determinism.}  
    The allocator deals in large aligned blocks. This keeps the early
    memory layout deterministic across boots, which greatly simplifies
    debugging: numerical addresses for key structures are stable or at
    least constrained.

  \item \textbf{Clear separation of roles.}  
    Lstack is strictly about \emph{physical} memory partitioning. It
    does not know about virtual addresses, page tables, or the notion of
    “heap”. This keeps it focused and makes it reusable for multiple
    purposes (heap base, reserved machine area, etc.).
\end{enumerate}

\subsection{libc / \texorpdfstring{\texttt{libc++}}{libc++} implications}

At this stage, \texttt{libc} is still not running, but the Lstack allocator
sets the stage for it:

\begin{itemize}
  \item The region carved out by Lstack will be split into:
    \begin{itemize}
      \item the \texttt{brk}-backed heap region, and
      \item the \texttt{mmap}-backed buddy-allocator region.
    \end{itemize}
  \item The alignment guarantees enforced here (64-byte chunks, later
        page-sized, etc.) propagate directly into the heap and buddy
        allocator invariants.
\end{itemize}

Instead of designing a heap allocator that can handle arbitrary,
fragmented physical memory, the kernel first re-shapes everything into a
nice big block and then builds a more sophisticated allocator on top of
that.

\section{Main Memory Selection and Reserved Machine Area}

\subsection{What happens}

During \texttt{Machine::init()}, the kernel chooses “its” main memory
region from what Lstack can provide and reserves a small portion for
internal use. This is implemented in the platform-specific machine code
(beyond what I extracted here) and exposed via the HAL.

From the trace:

\begin{verbatim}
→ os::Machine::init()
    → os::detail::Machine::init()
        → Machine::memory()
            → Lstack::allocate_largest()
        → INFO("[ Machine ] Main memory detected as 128839488 b")
        → INFO("[ Machine ] Reserving 1048576 b for machine use")
\end{verbatim}

The boot log matches this:

\begin{verbatim}
[ Machine ] Initializing heap
[ Machine ] Main memory detected as 128839488 b
[ Machine ] Reserving 1048576 b for machine use
\end{verbatim}

Conceptually, this step partitions physical memory into:

\begin{itemize}
  \item a large “main memory” block (\(\approx 123\)~MiB in my run),
  \item a 1~MiB chunk reserved for “machine use” (page tables, APIC
        structures, various kernel internal state).
\end{itemize}

\subsection{Why this split exists}

The separation between a “main memory” block and a “reserved machine
area” is motivated by:

\begin{enumerate}
  \item \textbf{Avoiding allocator self-interference.}  
    Some kernel structures (page tables, APIC mappings, low-level
    descriptors) need to be allocated before the general-purpose heap is
    ready, or must live outside it. Reserving a fixed-size region
    simplifies those early allocations.

  \item \textbf{Keeping the heap contiguous.}  
    By carving out the machine area as a single contiguous block, the
    remaining “main memory” region can be treated as one contiguous heap
    candidate. This improves locality and makes the buddy allocator
    implementation straightforward.

  \item \textbf{Predictable upper bound for kernel internals.}  
    A static 1~MiB reservation creates an implicit budget: kernel
    internals \emph{should} fit into that area. If they do not, some
    other change (e.g.\ more page tables or CPU features) is probably
    responsible, and that shows up early as a capacity issue.
\end{enumerate}

\subsection{libc / \texorpdfstring{\texttt{libc++}}{libc++} implications}

From \texttt{libc}'s point of view, the key outcomes of this phase are:

\begin{itemize}
  \item there is now a well-defined, contiguous physical region that
        will become the heap and \texttt{mmap} backend,
  \item kernel book-keeping is expected to use the reserved machine area,
        avoiding silent encroachment on the heap.
\end{itemize}

This explicit partitioning is particularly useful in a unikernel where
kernel and application share the same address space: it avoids subtle
interference where kernel growth eats into the application's heap
without leaving a clear trace.

\section{The Role of \texorpdfstring{\texttt{libc}}{libc} and \texorpdfstring{\texttt{libc++}}{libc++} in the Boot Process}

\subsection{Why \texttt{libc} exists in a C++ unikernel}

IncludeOS is implemented almost entirely in C++, and could in principle
run without a C library at all. \texttt{libc++} technically only relies
on a subset of the C library:

\begin{itemize}
  \item memory allocation (e.g.\ \texttt{malloc}/\texttt{free}),
  \item low-level memory operations (\texttt{memcpy}, \texttt{memset}, \dots),
  \item thread-local state (for \texttt{errno} and some C APIs),
  \item system time functions used indirectly by \texttt{std::chrono}.
\end{itemize}

A bespoke C++-only runtime could provide these directly. However,
IncludeOS intentionally ships and exposes a full musl \texttt{libc}. The
motivation is compatibility: third-party libraries overwhelmingly assume
the presence of a POSIX-like C environment. If kernel-provided APIs
diverged from standard \texttt{libc} interfaces, the ecosystem of
libraries usable inside services would shrink dramatically.

Thus the kernel does not merely use \texttt{libc} internally: it exposes
a complete \texttt{libc} ABI (Application Binary Interface) to the
application and any libraries linked into the unikernel. From the point
of view of a typical C or C++ library, the IncludeOS service “looks
like” an ordinary userspace process, even though underneath the OS and
the application are one and the same program.

\subsection{Why \texttt{libc} must be initialised before \texttt{libc++}}

Modern C++ runtime features impose several constraints on the order in
which \texttt{libc} and \texttt{libc++} can be brought up:

\begin{enumerate}
  \item \textbf{Global constructors often allocate and synchronise.}  
    Many global objects in \texttt{libc++} allocate memory or rely on
    function-local \texttt{static} variables protected by mutexes.
    These require:
    \begin{itemize}
      \item a working heap,
      \item working atomic operations,
      % \item working TLS for hidden per-thread state.
    \end{itemize}

  \item \textbf{Exceptions depend on thread-local state.}  
    \texttt{libc++abi} uses thread-local state maintained via \texttt{libc}
    TLS machinery. Without a valid TLS base, any throw/catch would be
    undefined behaviour as per C++11. TODO: We don't support TLS directly, how do we ensure this works as expected?

  \item \textbf{Standard synchronisation depends on libc locks.}  
    \texttt{std::mutex}, iostreams, and certain lazy initialisations in
    \texttt{libc++} depend on underlying locks implemented by musl. These
    locks are not usable until musl’s own initialisation sequence has
    run.

  \item \textbf{Time and timing facilities rely on libc’s view of the
    system clock.}  
    Both musl and \texttt{libc++} use system time facilities to implement
    \texttt{std::chrono::steady\_clock}, sleeps, and timed waits. In
    IncludeOS these go through \texttt{libc}'s syscall shims into the
    kernel’s timer infrastructure, and they are not meaningful until
    libc’s internal clock state is established.
\end{enumerate}

For these reasons, IncludeOS must bring up musl’s bootstrap machinery
(\texttt{\_\_libc\_start\_main} and related functions) before any C++
global constructors run. This ordering is reflected in the trace and in
\texttt{src/platform/x86\_pc/init\_libc.cpp}:

\begin{verbatim}
→ x86::init_libc(magic, addr)
    → build argv[] and environment strings
    → build auxv (ELF auxiliary vector)
    → setup syscall MSRs (IA32_STAR, IA32_LSTAR)
    → kernel::state().allow_syscalls = true
    → __libc_start_main(kernel_main, argc, argv)
\end{verbatim}

\subsection{Why expose a complete \texttt{libc} surface}

Even if \texttt{libc++} required only a minimal subset of C APIs,
restricting IncludeOS to that subset would significantly limit the set
of compatible libraries. Typical microservice stacks (JSON, HTTP,
crypto, gRPC, etc.) assume:

\begin{itemize}
  \item POSIX-style file and socket APIs,
  \item the existence of \texttt{errno},
  \item standard time functions,
  \item thread-local C state.
\end{itemize}

IncludeOS therefore exposes most of musl’s \texttt{libc} surface and
implements a compatibility layer (in \texttt{src/musl/*.cpp}) that maps
these C-level expectations onto the unikernel’s event-driven kernel. The
boot sequence described in this document is precisely what makes it
possible to present this familiar programming environment on top of a
single-address-space kernel.

\subsection{Consequences for the boot order}

The outcome of these design decisions is that \textbf{libc must be
partially initialised very early}:

\begin{itemize}
  \item the heap must exist,
  \item the virtual memory layout must be stable,
  \item the ELF auxiliary vector (auxv) must be populated,
  \item the syscall entry MSRs must be programmed,
  \item musl’s startup must execute before the first C++ global
        constructor.
\end{itemize}

This explains why the boot sequence performs:

\begin{enumerate}
  \item memory map parsing,
  \item symbol relocation,
  \item \texttt{.bss} initialisation,
  \item Lstack setup,
  \item heap and buddy allocator initialisation,
  \item paging and virtual memory layout setup,
  \item crash-context and IDT setup,
  \item and only then: musl \texttt{libc} bootstrap.
\end{enumerate}

From that point onward, C++ global constructors can safely run, and the
rest of the kernel can initialise higher-level structures without
risking undefined behaviour from an uninitialised runtime.

\section{Virtual Memory and Paging Layout}

\subsection{What happens}

The architecture-specific paging code is in
\texttt{src/arch/x86\_64/paging.cpp}. The entry point is a weak function
\texttt{\_\_arch\_init\_paging()}:

\begin{verbatim}
__attribute__((weak))
void __arch_init_paging() {
  INFO("x86_64", "Initializing paging");
  auto default_fl = Flags::present | Flags::writable | Flags::huge
                    | Flags::no_exec;
  __pml4 = new Pml4(0);
  ...
  INFO2("* Supported page sizes: %s", ...);
  INFO2("* Adding 512 1GiB entries @ 0x0 -> 0x8000000000");
  INFO2("* Identity mapping 0 -> 0, size 16.000_GiB, flags 0x%x", ...);
  INFO2("* Marking page 0 as not present");
  INFO2("* Allowing execute on 0x201000 -> 0x37dca8");
  INFO2("* Passing page tables to CPU");
}
\end{verbatim}

The log output in my run matches this exactly:

\begin{verbatim}
[ x86_64 ] Initializing paging
  * Supported page sizes: 4.000_KiB, 2.000_MiB
  * Adding 512 1GiB entries @ 0x0 -> 0x8000000000
  * Identity mapping 0 -> 0, size 16.000_GiB, flags 0x83
  * Marking page 0 as not present
  * Allowing execute on 0x201000 -> 0x37dca8
  * Passing page tables to CPU
\end{verbatim}

After paging is initialised, the kernel computes and prints a virtual
memory map from \texttt{src/kernel/memmap.cpp}, showing how different
regions of physical memory have been classified:

\begin{verbatim}
[ Kernel ] Total memory detected as 96.098_MiB
           Assigning fixed memory ranges (Memory map)
           * ELF symbol table size 0, ignoring.
           * Assigning heap 0x502000 -> 0x6018fff
[ Kernel ] Virtual memory map
           0x0000001000 - 0x0000007fff  Page tables
           0x0000010000 - 0x000009d3ff  Stack
           0x000009fc00 - 0x000009ffff  Reserved (Multiboot)
           0x00000f0000 - 0x00000fffff  Reserved (Multiboot)
           0x0000201000 - 0x000037dfff  ELF .text
           0x0000502000 - 0x0006018fff  Dynamic memory
           ...
\end{verbatim}

\subsection{Why it matters}

This phase:

\begin{itemize}
  \item sets up a large, identity-mapped region (16~GiB in the current
        configuration) so the kernel can treat physical and virtual
        addresses as the same in early code,
  \item unmapped page~0 to catch NULL-pointer dereferences as page
        faults,
  \item ensures that only the kernel’s \texttt{.text} region is marked
        executable (the \texttt{Allowing execute on} log line),
  \item carves out well-defined regions for stack, Multiboot structures,
        code, heap, and reserved areas.
\end{itemize}

The printed map doubles as documentation of the chosen layout and as a
debug tool when hunting memory corruption or mapping bugs.

\subsection{libc / \texorpdfstring{\texttt{libc++}}{libc++} requirements}

\texttt{libc} and \texttt{libc++} assume a flat virtual address space, but
they need:

\begin{itemize}
  \item executable code pages for function pointers and virtual calls,
  \item a mapped and reasonably sized stack,
  \item a mapped and writable heap.
\end{itemize}

All of that is established here, by the combination of
\texttt{\_\_arch\_init\_paging} and the memmap logic in
\texttt{memmap.cpp}.

\section{Heap, \texorpdfstring{\texttt{brk}}{brk}, \texorpdfstring{\texttt{mmap}}{mmap} and Buddy Allocator}

\subsection{What happens}

With main memory selected and paging in place, the kernel initialises
its heap semantics. The relevant code is spread across:

\begin{itemize}
  \item \texttt{src/kernel/heap.cpp},
  \item \texttt{src/musl/brk.cpp},
  \item \texttt{src/musl/mmap.cpp}.
\end{itemize}

The general flow is:

\begin{verbatim}
→ kernel::init_heap(free_mem_begin = 0x502000,
                    memory_end     = 0x6018fff)
    → kernel::heap_begin()
    → bits::roundto(4096, free_mem_begin)

    → __init_brk(begin = heap_base, size = 1 MiB)
    → __init_mmap(addr_begin = heap_base + brk_region,
                  size       = remaining)
        → os::mem::buddy::Alloc<false>::create(...)
\end{verbatim}

The log confirms the resulting heap region:

\begin{verbatim}
* Assigning heap 0x502000 -> 0x6018fff
0x0000502000 - 0x0006018fff, 91.090_MiB Dynamic memory
\end{verbatim}

The \texttt{brk}-backed part is used by musl for the traditional
contiguous heap; larger allocations and \texttt{mmap}-like calls are
serviced by the buddy allocator implemented in
\texttt{api/util/alloc\_buddy.hpp} and its backing code.

\subsection{Why it matters}

Musl uses two main backends for allocation:

\begin{itemize}
  \item a \texttt{brk}-based contiguous heap region for small/linear
        growth,
  \item an \texttt{mmap}-backed region for larger, page-oriented
        allocations.
\end{itemize}

The buddy allocator provides efficient splitting and coalescing of page
ranges for the latter. This is the point where the abstract “main
memory” chosen earlier becomes a concrete, structured heap.

\subsection{libc / \texorpdfstring{\texttt{libc++}}{libc++} requirements}

By the end of this phase:

\begin{itemize}
  \item \texttt{malloc}/\texttt{free} and \texttt{new}/\texttt{delete} have a working
        backend,
  \item the heap is page-aligned and non-overlapping with code or
        stack,
  \item allocation failures reflect genuine exhaustion, not uninitialised
        backends.
\end{itemize}

This is a hard precondition for any global C++ constructor that
allocates memory.

\section{RNG and Entropy}

\subsection{What happens}

The kernel seeds its RNG (Random Number Generator) in
\texttt{src/kernel/rng.cpp}, using CPU cycle counters and a
Keccak-based permutation:

\begin{verbatim}
→ RNG::init()
    → rng_reseed_init(fallback_entropy, rounds = 1)
    → reseed_now()
        → os::cycles_since_boot()
            → os::Arch::cpu_cycles()
        → rng_absorb(...)
        → keccak_1600_p(...)
    → per_cpu_help<rng_state, 32>(array)
\end{verbatim}

\subsection{Why it matters}

This provides a source of randomness early in boot, before other entropy
sources (for example hardware devices) are available. It can be used
for non-cryptographic randomness and as a seed for more sophisticated
RNGs.

\subsection{libc / \texorpdfstring{\texttt{libc++}}{libc++} requirements}

\texttt{libc}'s \texttt{rand()} does not require strong entropy, but
higher-level code may. For the build at this commit:

\begin{quote}
  TODO: document how TLS/crypto libraries (e.g.\ Botan, OpenSSL, s2n)
  hook into this RNG or provide their own entropy sources.
\end{quote}

\section{Crash Contexts and IDT}

\subsection{What happens}

The kernel sets up crash contexts and the Interrupt Descriptor Table
(IDT). The IDT code lives in \texttt{src/platform/x86\_pc/idt.cpp} and
the crash context support in \texttt{api/kernel/crash\_context.hpp} and
its implementation.

From the trace and source:

\begin{verbatim}
→ __init_crash_contexts()
→ x86::idt_initialize_for_cpu(0)
    → x86_IDT::init()
        → memset(entry[], 0, sizeof(entry))
    → x86_IDT::set_exception_handler(vec, __cpu_except_N)
    → x86_IDT::set_handler(vec, unused_interrupt_handler)
\end{verbatim}

This populates exception vectors (divide-by-zero, page fault, etc.) with
dedicated handlers and hooks up default handlers for unused IRQ lines.

\subsection{Why it matters}

Without a valid IDT, any exception would cause a double or triple fault
and reset the machine. This step ensures that:

\begin{itemize}
  \item all critical exception vectors have proper handlers,
  \item unused IRQs are routed to a safe default handler,
  \item crash contexts can capture register state for debugging.
\end{itemize}

\subsection{libc / \texorpdfstring{\texttt{libc++}}{libc++} requirements}

While user code does not interact with the IDT directly, it relies on
the fact that faults lead to controlled crashes with useful diagnostics
rather than abrupt resets. Any use of stack unwinding or signal-like
behaviour will also depend on this infrastructure being present.

\section{Libc Bootstrap: Auxv, Syscall MSRs, \texorpdfstring{\texttt{\_\_libc\_start\_main}}{__libc\_start\_main}}

\subsection{What happens}

The platform-specific \texttt{x86::init\_libc} function in
\texttt{src/platform/x86\_pc/init\_libc.cpp} prepares the environment
for musl and then hands over to \texttt{\_\_libc\_start\_main}.

Its responsibilities are:

\begin{enumerate}
  \item \textbf{Build \texttt{argv} and minimal environment.}  

    It constructs a small \texttt{argv} array (service name plus a few
    environment variables like \texttt{"PATH=/"} and \texttt{"LC\_ALL=C"}).

  \item \textbf{Populate the ELF auxiliary vector (auxv).}  

    It then uses an internal \texttt{auxv\_t} helper to fill in entries
    such as:

    \begin{itemize}
      \item \texttt{AT\_PAGESZ = 4096},
      \item \texttt{AT\_CLKTCK = 100},
      \item \texttt{AT\_PHENT = ehdr->e\_phentsize},
      \item \texttt{AT\_PHDR = \&phdr[0]},
      \item \texttt{AT\_PHNUM = ehdr->e\_phnum},
      \item \texttt{AT\_BASE = nullptr},
      \item \texttt{AT\_ENTRY = \&kernel\_main},
      \item \texttt{AT\_HWCAP = 0} (no extra hardware capability flags are
            advertised here),
      \item UID/GID/EUID/EGID all set to zero,
      \item \texttt{AT\_SECURE = 1},
      \item \texttt{AT\_PLATFORM = "x86\_64"},
      \item supplemental randomness (via musl’s \texttt{AT\_RANDOM}).
    \end{itemize}

  \item \textbf{Program syscall-related MSRs.}  

    For x86\_64, it sets up IA32\_STAR and IA32\_LSTAR:

    \begin{verbatim}
    uint64_t star_kernel_cs = 8ull << 32;
    uint64_t star_user_cs   = 8ull << 48;
    uint64_t star = star_kernel_cs | star_user_cs;
    x86::CPU::write_msr(IA32_STAR,  star);
    x86::CPU::write_msr(IA32_LSTAR, (uintptr_t) &__syscall_entry);
    \end{verbatim}

    The actual syscall entry stub \texttt{\_\_syscall\_entry} is implemented in
    \texttt{src/arch/x86\_64/\_\_syscall\_entry.asm} and is responsible
    for saving registers and forwarding to the musl syscall layer.

  \item \textbf{Allow syscalls and start musl.}  

    Finally, it marks syscalls as allowed in the kernel state and calls:

    \begin{verbatim}
    kernel::state().allow_syscalls = true;
    __libc_start_main(kernel_main, argc, argv.data());
    \end{verbatim}
\end{enumerate}

\subsection{Why it matters}

The auxiliary vector tells musl key properties such as page size,
program headers, and entrypoint. IA32\_STAR and IA32\_LSTAR configure the
hardware syscall path so that musl’s syscall wrappers can trap into the
kernel in a controlled way.

\subsection{libc / \texorpdfstring{\texttt{libc++}}{libc++} requirements}

By the end of this stage, musl and \texttt{libc++} can rely on:

\begin{itemize}
  \item valid auxv contents (including page size, program headers, and
        entrypoint),
  \item a functioning heap backend from earlier steps,
  \item a working syscall entry path for musl’s internal syscalls.
\end{itemize}

Global C and C++ constructors can now safely execute; the system log
constructor is one of the earliest to run.

\section{System Log and stdout Integration}

\subsection{What happens}

As part of musl’s startup, C++ global constructors run. One of these is
the system log glue in \texttt{src/kernel/system\_log.cpp}, which hooks
the system log into the OS-level stdout handler list.

Conceptually:

\begin{verbatim}
→ system_log_gconstr()
    → delegate<...>(SystemLog::write)
    → os::add_stdout(delegate)
\end{verbatim}

\texttt{os::add\_stdout} appends the delegate to a small
\texttt{Fixed\_vector<os::print\_func, 8>} of print handlers in
\texttt{src/kernel/kernel.cpp}. Later kernel code and the service may
emit log lines via this mechanism.

\subsection{Why it matters}

Once global constructors and later kernel code start to log, there needs
to be a consistent place for that output to go. Doing this inside musl’s
initialisation path ensures it is in place before user-level service
code runs.

\subsection{libc / \texorpdfstring{\texttt{libc++}}{libc++} requirements}

\begin{itemize}
  \item \texttt{printf}, \texttt{puts}, and C++ streams can be bridged to
        these stdout handlers via musl and \texttt{libc++}'s iostream
        implementation.
  \item Logging from inside constructors, or from early utility code,
        must not crash because no stdout is registered.
\end{itemize}

\section{ACPI, SMP, SMBIOS Discovery and \texttt{\_\_platform\_init}}

\subsection{What happens}

After \texttt{\_\_libc\_start\_main} has started \texttt{kernel\_main} and
run C++ constructors, the platform-specific initialisation function
\texttt{\_\_platform\_init} in \texttt{src/platform/x86\_pc/platform.cpp}
is responsible for bringing up hardware-specific subsystems:

\begin{verbatim}
void __platform_init()
{
  // read ACPI tables
  x86::ACPI::init();

  // read SMBIOS tables
  x86::SMBIOS::init();

  INFO("x86", "Setting up GDT, TLS, IST");
#ifdef ARCH_x86_64
  x86::ist_initialize_for_cpu(0, 0x9D3F0);
#endif

  INFO("x86", "Initializing CPU 0");
  x86::initialize_cpu_tables_for_cpu(0);
  Events::get(0).init_local();

  // setup APIC, APIC timer, SMP etc.
  x86::APIC::init();

  // enable interrupts
  asm volatile("sti");

  // optional SMP init
  x86::init_SMP();  // if enabled

  INFO("x86", "Setting up kernel clock sources");
  x86::Clocks::init();
  ...
  x86::APIC_Timer::calibrate();

  INFO("Initializing drivers");
  kernel::run_ctors(&__driver_ctors_start, &__driver_ctors_end);

  hw::PCI_manager::init();
  hw::PCI_manager::init_devices(PCI::STORAGE);
  hw::PCI_manager::init_devices(PCI::NIC);
  os::machine().print_devices();
}
\end{verbatim}

The logs I see correspond to the first half of this sequence; the later
parts (driver ctors, PCI scanning) are not yet visible in the current
trace because the APIC timer calibration is still ongoing.

\subsection{ACPI: discovering LAPIC and IOAPIC}

ACPI support is implemented in \texttt{src/platform/x86\_pc/acpi.cpp} and
\texttt{acpi.hpp}. \texttt{x86::ACPI::init()} simply calls
\texttt{ACPI::discover()}, which:

\begin{enumerate}
  \item looks for the RSDP (Root System Description Pointer) at a QEMU
        “guess” address (\texttt{0xF68C0}),
  \item if not found there, scans the BIOS region \texttt{0xE0000} to
        \texttt{0xFFFFF} for the signature \texttt{"RSD PTR "},
  \item verifies checksums,
  \item calls \texttt{ACPI::begin(rsdp)} once a valid RSDP is found.
\end{enumerate}

\texttt{ACPI::begin} prints:

\begin{verbatim}
[ ACPI ] Reading headers
         OEM: BOCHS  Rev. 0
\end{verbatim}

and then walks the system description table (RSDT/XSDT), collecting:

\begin{itemize}
  \item LAPIC (Local APIC) entries,
  \item IOAPIC (I/O APIC) entries,
  \item interrupt source overrides (IRQ remapping),
  \item NMI (Non-Maskable Interrupt) configuration.
\end{itemize}

The APIC information printed in the log comes from this analysis:

\begin{verbatim}
[ ACPI ] Reading APIC information
         LAPIC base: 0xfee00000  (flags: 0x1)
         I/O APIC 0   ADDR 0xfec00000  INTR 0x0
         IRQ redirect for bus 0 from IRQ 0  to VEC 2
         IRQ redirect for bus 0 from IRQ 5  to VEC 5
         IRQ redirect for bus 0 from IRQ 9  to VEC 9
         IRQ redirect for bus 0 from IRQ 10 to VEC 10
         IRQ redirect for bus 0 from IRQ 11 to VEC 11
\end{verbatim}

\subsection{SMP discovery}

ACPI’s MADT (Multiple APIC Description Table) parsing also counts CPUs
and records LAPIC IDs. This is reported as:

\begin{verbatim}
[ SMP ] Found 1 APs
\end{verbatim}

The actual SMP bring-up is guarded by \texttt{INCLUDEOS\_SMP\_ENABLE}
and implemented in \texttt{src/platform/x86\_pc/smp.cpp}; in my current
configuration only CPU~0 is active.

\subsection{SMBIOS}

SMBIOS parsing is implemented in
\texttt{src/platform/x86\_pc/smbios.cpp}. \texttt{x86::SMBIOS::init()}
locates the SMBIOS entry point and walks the tables, logging:

\begin{verbatim}
[ SMBIOS ] Version 2.8
           BIOS vendor: rel-1.16.3-0-ga6ed...
           Manufacturer: Standard PC (i440FX + PIIX, 1996)
           Product name: pc-i440fx-9.2
           Physical memory array with 128 MB capacity
\end{verbatim}

This is primarily used for information and for potential future
platform-specific behaviour.

\subsection{Why it matters}

\begin{itemize}
  \item ACPI identifies LAPIC and IOAPIC addresses and IRQ routing, which
        are crucial for APIC initialisation.
  \item SMP information indicates how many application processors are
        present and what their APIC IDs are.
  \item SMBIOS provides a standard view of the platform (vendor, model,
        memory configuration) useful for logging and potential quirks.
\end{itemize}

\subsection{libc / \texorpdfstring{\texttt{libc++}}{libc++} requirements}

These steps are mainly OS/hardware concerns and do not directly affect
\texttt{libc}/\texttt{libc++}. However, any multithreading
(\texttt{std::thread}, POSIX threads) will be constrained by whatever
SMP setup is established here.

\section{GDT, TLS, IST and CPU 0 Initialisation}

\subsection{What happens}

The same \texttt{\_\_platform\_init} function continues by setting up
long-mode segmentation and per-CPU state:

\begin{verbatim}
INFO("x86", "Setting up GDT, TLS, IST");
#ifdef ARCH_x86_64
  x86::ist_initialize_for_cpu(0, 0x9D3F0);
#endif

INFO("x86", "Initializing CPU 0");
x86::initialize_cpu_tables_for_cpu(0);
Events::get(0).init_local();
\end{verbatim}

The IST (Interrupt Stack Table) logic is implemented in
\texttt{src/arch/x86\_64/ist.cpp}. \texttt{ist\_initialize\_for\_cpu}
allocates separate stacks for:

\begin{itemize}
  \item regular interrupts,
  \item NMIs,
  \item double-fault handling.
\end{itemize}

It stores stack pointers in the AMD64 Task State Segment (TSS) and
ensures they are aligned, writable and tested (by forcing a write near
the top of each stack).

\texttt{x86::initialize\_cpu\_tables\_for\_cpu} (implemented in
\texttt{src/platform/x86\_pc/gdt.cpp} and \texttt{idt.cpp}) then:

\begin{itemize}
  \item sets up the final GDT entries (code, data, TSS),
  \item loads the TSS via \texttt{ltr},
  \item associates the IDT prepared earlier with CPU~0,
  \item initialises the \texttt{Events} structure for the local CPU.
\end{itemize}

\subsection{Why it matters}

\begin{itemize}
  \item The GDT defines code/data segments and the TSS used by the CPU
        for privilege transitions and IST.
  \item IST entries provide dedicated stacks for critical exceptions
        (NMI, double-fault) that must not reuse the normal kernel stack.
  \item \texttt{Events::get(0).init\_local()} prepares the per-CPU event
        queue used by the IncludeOS event-driven model.
\end{itemize}

\subsection{libc / \texorpdfstring{\texttt{libc++}}{libc++} requirements}

Long mode and basic TLS were already functional before this point, but
this phase strengthens those guarantees and ensures exceptions use safe
stacks. Any subsequent fault or interrupt handling by the kernel occurs
with a well-defined stack and descriptor state.

\section{IOAPIC and xAPIC Setup, Enabling Interrupts}

\subsection{What happens}

Next, the kernel brings IOAPIC and LAPIC online and enables interrupts.
This logic lives mainly in:

\begin{itemize}
  \item \texttt{src/platform/x86\_pc/apic.cpp},
  \item \texttt{src/platform/x86\_pc/ioapic.cpp},
  \item \texttt{src/platform/x86\_pc/apic\_timer.cpp}.
\end{itemize}

From \texttt{\_\_platform\_init} and the logs:

\begin{verbatim}
[ I/O APIC ] Initializing ID=0
             Addr: 0xfec00000  Version: 0x20  Intr: 0  Entries: 24
[ xAPIC ] Enabling xAPIC
           ID: 0  Ver: 50014
[ x86 ] Enabling interrupts
\end{verbatim}

In code:

\begin{verbatim}
x86::APIC::init();
asm volatile("sti");
\end{verbatim}

\texttt{x86::APIC::init} configures the local APIC at the physical base
reported by ACPI (\texttt{0xFEE00000}) and initialises the IOAPIC at
\texttt{0xFEC00000}, programming redirection entries according to ACPI
overrides.

After \texttt{sti}, interrupts begin to flow: timer interrupts, any
device IRQs that have been unmasked, and IPIs (in SMP configurations).

\subsection{Why it matters}

Without a working IOAPIC/LAPIC configuration, the kernel cannot:

\begin{itemize}
  \item handle timer interrupts,
  \item respond to device interrupts,
  \item drive an event loop.
\end{itemize}

This is the first point where the system becomes truly interrupt-driven.

\subsection{libc / \texorpdfstring{\texttt{libc++}}{libc++} requirements}

Once interrupts are enabled, \texttt{libc} operations may interleave with
IRQ handlers. This means:

\begin{itemize}
  \item any shared state used in both contexts must be properly locked,
  \item time-related syscalls can start to rely on clock sources once
        calibration is complete.
\end{itemize}

\section{CMOS Clock, CPU Frequency, and APIC Timer Calibration}

\subsection{What happens}

Timekeeping and CPU frequency estimation are set up in:

\begin{itemize}
  \item \texttt{src/platform/x86\_pc/cmos.cpp},
  \item \texttt{src/platform/x86\_pc/cmos\_clock.cpp},
  \item \texttt{src/platform/x86\_pc/clocks.cpp},
  \item \texttt{src/platform/x86\_pc/apic\_timer.cpp}.
\end{itemize}

The log shows:

\begin{verbatim}
[ x86 ] Setting up kernel clock sources
[ CMOS ] Enabling regular clock sync for CMOS clock
[ x86 ] CMOS clock initialized
[ CMOS ] Estimating CPU-frequency
         |
         +--(18 samples, 0.000100 sec. interval)
         |
         Enabled redirected entry 2 ioapic 0 -> 0 on apic 0
         +--> 3819.505224 MHz
[ APIC ] Measuring APIC timer...
\end{verbatim}

\texttt{x86::Clocks::init()} configures the CMOS-based wall clock and a
TSC (Time Stamp Counter) based monotonic clock. \texttt{CMOS\_clock} is
used to estimate CPU frequency by sampling TSC deltas over a fixed
interval (18 samples at 100~µs in my build).

The APIC timer calibration in \texttt{x86::APIC\_Timer::calibrate()} then
programs the local APIC timer and measures how many APIC ticks correspond
to a given real-time duration, using the newly estimated CPU frequency
as a reference.

While this is happening, the APIC timer fires periodically and the
xAPIC interrupt handler in \texttt{src/platform/x86\_pc/apic\_asm.asm}
and \texttt{apic.cpp} is hit repeatedly. This is where my captured logs
stop; the system is running and handling APIC interrupts but has not yet
completed calibration.

\subsection{Why it matters}

\begin{itemize}
  \item CPU frequency is needed to translate cycle counts (from the TSC)
        into time units.
  \item APIC timer calibration is needed to translate APIC ticks into
        real-time durations.
\end{itemize}

Without these calibrations, time-based functionality (timeouts, sleeps,
scheduling) would be inaccurate or unavailable.

\subsection{libc / \texorpdfstring{\texttt{libc++}}{libc++} requirements}

Functions like \texttt{clock\_gettime}, \texttt{sleep},
\texttt{nanosleep}, and the \texttt{std::chrono} APIs rely on calibrated
clocks. In this build, musl’s time syscalls are implemented via shims in
\texttt{src/musl/clock\_gettime.cpp}, \texttt{gettimeofday.cpp},
\texttt{nanosleep.cpp}, etc., which ultimately call into the kernel’s
clock and timer infrastructure initialised here.

\begin{quote}
  TODO: explicitly map each musl time syscall in \texttt{src/musl/*.cpp}
  to the underlying kernel API calls, and document their precision and
  limitations.
\end{quote}

\section{What This Trace Does Not Cover (Yet)}

The GDB trace and log stop during APIC timer calibration. In this trace,
there is no direct evidence for:

\begin{itemize}
  \item the exact runtime ordering of:
    \begin{itemize}
      \item \texttt{kernel::start()},
      \item \texttt{os::start()},
      \item \texttt{Service::start()} or \texttt{kernel\_main()},
    \end{itemize}
  \item PCI scanning and driver initialisation (virtio, vmxnet3, \dots),
    although \texttt{\_\_platform\_init} clearly calls
    \texttt{hw::PCI\_manager::init()} and runs driver constructors
    after APIC timer calibration,
  \item network stack bring-up (ARP, IP, TCP/UDP),
  \item entering the main event loop (\texttt{os::event\_loop()}),
  \item user/service-level code running.
\end{itemize}

From the source (\texttt{src/platform/x86\_pc/platform.cpp}) we know
that once APIC timer calibration has completed, driver constructors and
PCI probing run, followed by storage and NIC device initialisation.
However, the current logs do not reach those points, so the timing of
service-level entry relative to these operations is not yet empirically
verified.

Next steps for this document will either:

\begin{itemize}
  \item extend the trace to cover these later stages, or
  \item inspect the relevant source files (\texttt{src/kernel/os.cpp},
        the service entry points, driver code) and explicitly mark those
        sections as “derived from code” rather than “derived from logs”.
\end{itemize}

\section{Summary of libc / \texorpdfstring{\texttt{libc++}}{libc++} Preconditions by Phase}

For convenience, this section summarises which invariants are
established before \texttt{libc}/\texttt{libc++} rely on them:

\begin{enumerate}
  \item After \texttt{.bss} clear:
    \begin{itemize}
      \item globals are zeroed.
    \end{itemize}

  \item After Machine + Lstack + heap initialisation:
    \begin{itemize}
      \item page-aligned heap and \texttt{mmap} regions exist,
      \item \texttt{malloc}/\texttt{new} have valid backends.
    \end{itemize}

  \item After paging layout:
    \begin{itemize}
      \item code, stack, and heap are mapped and non-overlapping,
      \item page~0 is unmapped,
      \item only \texttt{.text} is executable.
    \end{itemize}

  \item After RNG:
    \begin{itemize}
      \item a basic kernel RNG is seeded.
    \end{itemize}

  \item After crash contexts + IDT:
    \begin{itemize}
      \item exceptions generate crash contexts rather than resets.
    \end{itemize}

  \item After \texttt{x86::init\_libc}:
    \begin{itemize}
      \item auxv is populated,
      \item syscall MSRs (IA32\_STAR, IA32\_LSTAR) are configured,
      \item \texttt{\_\_libc\_start\_main} is running.
    \end{itemize}

  \item After system log constructors:
    \begin{itemize}
      \item logging backends are registered via \texttt{os::add\_stdout}.
    \end{itemize}

  \item After ACPI/IOAPIC/xAPIC and clock setup:
    \begin{itemize}
      \item interrupts are enabled,
      \item CPU frequency and APIC timer calibration are underway.
    \end{itemize}
\end{enumerate}

By the end of the observed trace, \texttt{libc} and \texttt{libc++} have
the basic infrastructure they need, but the kernel is still refining its
timekeeping and interrupt behaviour. Service-level code and the
high-level OS facilities come afterwards.

\end{document}
